\begin{definition}[Lecteur de signature et transport des représentants]
\label{def:firma555}
Soit \(\mathcal X_\times=D_9^2\times\{N,E,S,O\}\), de cardinal \(324\).
L'avancée \(P\) déplace d'une case dans la direction mémorisée et le demi-
tour \(H\) inverse cette direction. Le lecteur utilise
\[
\ell_9(1,2,4,8,7,5)=(0,1,2,3,4,5),\qquad
\ell_9(3)=\ell_9(6)=\ell_9(9)=0.
\]
Dans chacune des six fenêtres de neuf étapes, on évalue
\(\ell_9(\rho_9(ij))\) avant l'avancée produit sélectionnée par TRIT aux
temps \(2,5,8\pmod9\). La somme de chaque fenêtre est réduite modulo
\(10\). Après les étapes \(27\) et \(54\), la direction est inversée.
La règle détermine \(E_{\rm sig}(x)\) à partir de la position et de la direction
de \(x\), avec les conventions temporelles du noyau.
\end{definition}

\begin{proposition}[Raffinement minimal stable de la signature]
\label{prop:firma555}
En partant de l'équivalence par signature \(\sim_0\), l'itération
\[
x\sim_{n+1}y\ \Longleftrightarrow\
x\sim_n y,\quad Px\sim_n Py,\quad Hx\sim_n Hy
\]
se termine et produit la congruence stable la plus grossière qui raffine la signature. Dans le domaine précédent, le recensement est \(26\to28\to28\) ; la fibre de \(555555\) se sépare en trois classes de huit représentants.
\end{proposition}
\begin{proof}
À chaque étape, on raffine une partition d'un ensemble fini. Le processus se
stabilise après un nombre fini de scissions. Au point fixe,
l'équivalence est préservée sous \(P\) et \(H\). Si une congruence stable
raffine \(\sim_0\), elle raffine par récurrence chaque \(\sim_n\) : ses images
par les deux opérateurs restent équivalentes. Elle raffine donc le point
fixe obtenu, ce qui prouve sa minimalité en information ajoutée.

L'énumération finie applique la définition aux \(81\) positions et aux quatre orientations. On regroupe d'abord leurs six sommes modulo \(10\), puis le triplet formé par la classe de \(x\), celle de \(Px\) et celle de \(Hx\), jusqu'à stabilisation. L'histogramme résultant est \(27\cdot8+108=324\), et la seule scission de signature est celle indiquée. Les représentants suivants explicitent les trois images :
\[
\begin{array}{c|c|c}
x&E_{\rm sig}(x)&E_{\rm sig}(Px)\\\hline
(3,5,N)&555555&555177\\
(3,4,N)&555555&555717\\
(3,4,S)&555555&555771
\end{array}
\]
La règle complète et son énumération exécutable sont conservées dans le paquet
de reproduction ; les témoins du tableau montrent en outre que conserver
seulement la signature ne rend pas sa continuation univoque.
\end{proof}

La conclusion décrit l'information minimale du transport, et non une
nouvelle masse. Le facteur additif conserve dix-huit représentants par signature
(neuf positions compatibles et deux orientations). En le combinant avec les
fibres produit de tailles \(24\) et \(8\), on obtient respectivement
\(432\) et \(144\). L'incidence régionale et la mémoire de feuillet conservent
ainsi des fonctions distinctes dans le même catalogue.
